\documentclass[11pt]{article}
\usepackage[margin=1in]{geometry}
\usepackage{parskip}
\usepackage[T1]{fontenc}
\usepackage[utf8]{inputenc}
\usepackage{times}
\begin{document}
\thispagestyle{empty}

\noindent Ku Kang\\
Chemical, Biological, Radiological and Nuclear (CBRN) Defense Research Institute\\
Seoul, Republic of Korea\\
bisu9082@gmail.com\\

\vspace{0.5em}
\noindent 8 September 2026

\vspace{1em}
\noindent Dear Editors,

We are pleased to submit our manuscript, ``\textbf{How much of hazard news is news? Decomposing open-source coverage of chemical accidents and CBRN threats for disaster risk reduction},'' for consideration as a Research Article in the \emph{International Journal of Disaster Risk Reduction}.

Emergency planners and threat analysts increasingly read the volume of news coverage as a proxy for how often a hazard occurs, especially for chemical, biological, radiological and nuclear (CBRN) threats, for which no incident registry exists. The manuscript tests the assumption behind that practice: that an article about a hazard is evidence that something happened. Using 109,610 Korean news articles from 54 outlets over five years, for 30 queries in three blocks (industrial chemical accidents, CBRN security threats, conventional terrorism), we decompose news volume into five labelled components and show that only one article in three reports an event, one in ten for CBRN security hazards by model coding, and one in fifty a domestic CBRN event. Two of the authors independently coded a 400-article stratified subset (inter-coder $\kappa = 0.85$) and put the CBRN figure lower still, at one article in twenty (5.0\%, 95\% CI 2.5--7.9), and the domestic share at one in forty, because the model coders counted topically adjacent occurrences that the codebook excludes. We report which results that validation licenses and which remain on model labels. Query-level event shares, computed on the model labels, separate the CBRN block from both control blocks, which this design lacks the power to distinguish from each other; the separation survives conditioning on topical relevance, so it is not produced by the three query sets retrieving with different precision. We define a reference-bias metric (the share of articles citing a past anchor event), which reaches 89\% for Novichok and 33\% for hydrogen fluoride against 1--2\% for matched controls, and show against the national chemical-accident registry that article counts for six substances are unrelated to incident counts. A lightweight text classifier recovers the event component, and the archive's own incident taxonomy does not improve it. The human validation is also a result in its own right: two language models agreed with each other at $\kappa = 0.79$ while agreeing with human coders at 0.50, and at 0.24 in the CBRN block, which is a caution for the growing use of model-based annotation in this literature.

We believe the work fits the journal's scope for three reasons. It gives practitioners an operational quantity, the event share with a confidence interval, for a widely used but untested proxy; it benchmarks open-source coverage against an official disaster registry, which to our knowledge has not been done for chemical accidents; and it names and measures a form of media bias (reference bias) that the selection- and description-bias literature does not capture. The three-block design with a registry-backed control lets us state that CBRN coverage is unusual, not merely inflated, and a sensitivity analysis shows that the contrast survives the exclusion of every query a reader might assign differently. The paper closes with the procedure an agency can run on its own hazard and archive, and the codebook, anchor dictionary, labelled sample and classifier are released so that it can be applied elsewhere.

The manuscript has not been published previously and is not under consideration by another journal. All authors have approved the submission. We declare no competing interests. Two large language models were used as independent coders under a written codebook, and language models assisted with drafting; both uses are stated in the manuscript, and the authors take full responsibility for the content. Article identifiers, the gold-set labels including both human coders' labels, the codebook, the registry extract and all analysis code are already public at https://github.com/bisu9082/refbias.

Thank you for considering our manuscript.

\vspace{1em}
\noindent Sincerely,

\vspace{1.5em}
\noindent Ku Kang, on behalf of all authors\\
Miae Kim, Janghyeon Cha, Myeongsik Shin, Gyungjun Choi, Ku Kang

\end{document}
